\section{具有性质（GDF）的几类空间}

我们已经知道每个 Fréchet 空间都具有性质（GDF）（TVS, I, §3, No.~3, 定理 1 的推论 5）。

命题 2. --- 设 E 是一个向量空间， \((\mathrm{F}_{\alpha})_{\alpha\in\mathrm{A}}\) 是具有性质（GDF）的局部凸空间的一个族，且对每个 \(\alpha\in A\) ，设 \(h_{\alpha}\) 是 \(F_{\alpha}\) 到 E 中的一个线性映射。若 E 装备使诸 \(h_{\alpha}\) 连续的最细局部凸拓扑，则 E 具有性质（GDF）。

设 \(u\) 是 \(\mathbf{E}\) 到一个 Banach 空间 \(\mathbf{B}\) 中的线性映射，使得在 \(\mathbf{E} \times \mathbf{B}\) 中，图 \(\Gamma\) （ \(u\) 的图）的点的每个收敛序列的每个极限也属于 \(\Gamma\) 。只需证明对每个 \(\alpha \in \mathbf{A}\) ， \(u \circ h_{\alpha}\) 在 \(\mathrm{F}_{\alpha}\) 上连续（TVS, II, §4, No.~4, 命题 5）。现在，设 \((x_n)\) 是 \(\mathrm{F}_{\alpha}\) 的元素的一个序列，具有极限 \(a\) ，且使得序列 \(\left(u(h_{\alpha}(x_n))\right)\) 具有一个极限 \(b \in \mathbf{B}\) 。由于 \(h_{\alpha}\) 连续， \(h_{\alpha}(a)\) 是序列 \((h_{\alpha}(x_n))\) 在 \(\mathbf{E}\) 中的一个极限；因此由假设 \(b = u(h_{\alpha}(a))\) ，且由于 \(\mathrm{F}_{\alpha}\) 具有性质（GDF）， \(u \circ h_{\alpha}\) 连续。

推论. --- 具有性质（GDF）的局部凸空间的每个商空间都具有性质（GDF）。

命题 3. --- 自反 Fréchet 空间的强对偶具有性质（GDF）。

这是命题 2 与下述引理（或 TVS, IV, §3, No.~4, 命题 4）的推论：

引理. --- 设 F 是一个 Fréchet 空间， \(F'\) 是它的强对偶， \(F''\) 是它的双对偶。若 \(F''\) 的每个关于 \(\sigma(F'', F')\) 有界的子集都含于 F 的一个有界子集的闭包（关于 \(\sigma(F'', F')\) ）中，则 \(F'\) 是 Banach 空间的一个序列的直接极限。

因为，设 \((\mathrm{V}_{n})\) 是 F 中 0 的凸、均衡、闭邻域的一个递减的基本序列。对每个整数 n ，设 \(G_{n}\) 是 \(F'\) 的线性子空间，由极集 \(V_{n}^{\circ}\) （ \(V_{n}\) 的极集）生成。在 \(G_{n}\) 中， \(V_{n}^{\circ}\) 是一个吸收性凸集，因此它的度规 \(p_{n}\) 是 \(G_{n}\) 上的一个范数；此外， \(V_{n}^{\circ}\) 是强对偶 \(F'\) 的完备子集（TVS, III, §3, No.~8, 命题 11）；从而 \(G_{n}\) （装备范数 \(p_{n}\) ）是一个 Banach 空间（GT, III, §3, No.~5, 命题 9 的推论 2）。我们将证明 \(F'\) 上的强拓扑是诸 \(G_{n}\) 上这些 Banach 空间拓扑的直接极限，或者等价地说，一个强闭、均衡、凸子集 U （ \(F'\) 的子集）是 0 的强邻域，当且仅当它吸收每个 \(V_{n}^{\circ}\) 。这个条件是必要的，这是显然的；为看出它是充分的，只需证明 U 含有 \(F'\) 的一个桶。事实上，此时它的极集 \(U^{\circ}\) （在 \(F''\) 中）将关于 \(\sigma(\mathrm{F}''\mathrm{,}\mathrm{F}')\) 有界，因而由假设含于 F 的某个有界子集 B 的闭包（关于 \(\sigma(\mathrm{F}''\mathrm{,}\mathrm{F}')\) ）中，由此可以断定 U （它关于 \(\sigma(\mathrm{F}^{\prime},\mathrm{F}''\mathrm{)})\) 是闭的）含有 0 的强邻域 \(B^{\circ}\) （TVS, II, §6, No.~3, 定理 1 与 §8, No.~4）。

由假设，对每个整数 n 存在一个数 \(\lambda_{n}>0\) ，使得 \(\lambda_{n}V_{n}^{\circ}\subset\frac{1}{2}U\) ；设 \(A_{n}\) 是诸 \(\lambda_{i}V_{i}^{\circ}\) （ \(i\leqslant n\) ）的并的凸包。则对每个 n 有 \(A_{n}\subset\frac{1}{2}U\) ；设 W 是诸 \(A_{n}\) 的并； W 是含于 \(\frac{1}{2}U\) 中的一个凸、均衡、吸收集，只需证明它的强闭包（它是一个桶）含于 U 中。

于是，设 \(x'\) 是 \(F'\) 中不属于 U 的一个点。由于每个 \(V_n^\circ\) 关于 \(\sigma(F', F)\) 是紧的，每个 \(A_n\) 亦然（TVS, II, §2, No.~6, 命题 15），且由于 \(x' \notin 2A_n\) ，存在一个元素 \(x_n\) ，它属于 \(A_n\) 在 F 中的极集，使得 \(\langle x', x_n \rangle = 2\) （TVS, II, §5, No.~3, 命题 4）。序列 \((x_n)\) 在 F 中有界：因为每个 \(y' \in F'\) 都属于某个 \(V_k^\circ\) ，于是 \(|\langle y', x_n \rangle| \leqslant \lambda_k^{-1}\) （对 \(n \geqslant k\) ），由此得我们的断言（TVS, IV, §1, No.~1, 命题 1）。设 C 是 F 中包含所有 \(x_n\) 的一个有界、均衡凸集；则 \(C^\circ\) 是 \(F'\) 中 0 的一个邻域，而极集 \(C^{\circ\circ}\) （ \(C^\circ\) 在 \(F''\) 中的极集）关于 \(\sigma(F'', F')\) 是紧的（TVS, III, §3, No.~4, 命题 4 的推论 2 与 No.~5, 命题 7）。于是可以看出，序列 \((x_n)\) 有一个聚点 \(x''\) （在 \(F''\) 中，关于 \(\sigma(F'', F')\) ）；显然 \(\langle x', x'' \rangle = 2\) ，而另一方面， \(x''\) 属于 \(A_n\) 在 \(F''\) 中的极集（对每个 \(n\) ），从而属于极集 \(W^\circ\) （ W 在 \(F''\) 中的极集）。由此得出 \(x' \notin W^{\circ\circ}\) ，因而 x' 不在 W 关于 \(\sigma(F', F'')\) 的闭包中（TVS, II, §6, No.~3, 定理 1 与 §8, No.~4），从而更不在强拓扑的闭包中，证明至此完成。
% semantic-fragments: legacy-input-seam
\relax{}
